
\section*{8. The Case of a Base Field}
\phantomsection
\label{sga3:VIIA:8}
\addcontentsline{toc}{section}{8. The Case of a Base Field}

\medskip
\noindent\textbf{8.1.}
\phantomsection
\label{sga3:VIIA:8.1}

We now summarize the results obtained when \(S\) is the spectrum
of a field \(k\) of characteristic \(p>0\). We call a
\(k\)-group scheme \emph{algebraic} if its underlying scheme is
of finite type over \(k\). \SGARefLink{sga3:VIIA:7.2}{Theorem 7.2} then gives the following
result.\footnote{The editors supplied the numbering \SGARefLink{sga3:VIIA:8.1.1}{8.1.1}
through \SGARefLink{sga3:VIIA:8.1.3}{8.1.3} in order to highlight the results stated here.}

\medskip
\noindent\textbf{Theorem 8.1.1.}
\phantomsection
\label{sga3:VIIA:8.1.1}
\textit{The functor \(\mathfrak G_p\), which associates with
every finite-dimensional restricted \(p\)-Lie algebra
\(\mathcal L\) over \(k\) the \(k\)-group
\(\mathfrak G_p(\mathcal L)\), is left adjoint to the functor
which associates \(\operatorname{Lie}(G)\) with every algebraic
\(k\)-group \(G\).}

Combining this with \SGARefLink{sga3:VIIA:7.4}{Theorem 7.4} gives:

\medskip
\noindent\textbf{Theorem 8.1.2.}
\phantomsection
\label{sga3:VIIA:8.1.2}
\textit{The functors
\[
  \mathfrak G_p:
  \mathcal L\longmapsto\mathfrak G_p(\mathcal L),
  \qquad
  G\longmapsto\operatorname{Lie}(G),
\]
induce quasi-inverse equivalences between the category of
finite-dimensional restricted \(p\)-Lie algebras over \(k\)
and the category of algebraic \(k\)-groups of height at most
\(1\).}

Since \(\mathfrak G_p\) is a left adjoint, it commutes with
inductive limits,\footnote{The editors expanded the argument
which follows, as well as the proof of \SGARefLink{sga3:VIIA:8.1.3}{Corollary 8.1.3}.} and in
particular with cokernels. Let
\[
  \varphi:G\longrightarrow H,
  \qquad
  \varphi':G'\longrightarrow H
\]
be two morphisms between algebraic \(k\)-groups of height at
most \(1\). The fiber product \(G\times_HG'\) is again an
algebraic \(k\)-group of height at most \(1\), because the
Frobenius morphism
\[
  \operatorname{Fr}:G\longrightarrow G^{(p)}
\]
commutes with fiber products. Thus the inclusion of the
category of algebraic \(k\)-groups of height at most \(1\) into
the category of all algebraic \(k\)-groups commutes with fiber
products, and hence in particular with kernels. We obtain:

\medskip
\noindent\textbf{Corollary 8.1.3.}
\phantomsection
\label{sga3:VIIA:8.1.3}
\textit{The functor \(\mathfrak G_p\) is exact in the following
sense. Let
\[
  \pi:\mathcal L_1\longrightarrow\mathcal L_2
\]
be a surjective homomorphism between finite-dimensional
restricted \(p\)-Lie algebras over \(k\), and let \(i\) be the
inclusion of \(\mathcal L_0=\ker\pi\) into \(\mathcal L_1\).
Then the following sequence of algebraic \(k\)-groups is exact:}
\begin{center}
  \phantomsection
  \label{sga3:VIIA:8.1.3-Gp-exact-sequence}
  \(\displaystyle
1 \longrightarrow \mathfrak G_p(\mathcal L_0)
  \xrightarrow{\mathfrak G_p(i)}
\mathfrak G_p(\mathcal L_1)
  \xrightarrow{\mathfrak G_p(\pi)}
\mathfrak G_p(\mathcal L_2)
  \longrightarrow 1 .
\)

\end{center}
\footnote{Moreover, by VIA, \SGARefLink{thm:sga3-VIA-3.2}{3.2},
\(\mathfrak G_p(\mathcal L_2)\) represents the (fppf) quotient
sheaf of \(\mathfrak G_p(\mathcal L_1)\) by
\(\mathfrak G_p(\mathcal L_0)\).}

Indeed, by the preceding discussion, \(\mathfrak G_p(i)\)
induces an isomorphism from
\(\mathfrak G_p(\mathcal L_0)\) onto
\(\ker(\mathfrak G_p(\pi))\). This kernel is the same whether
it is taken in the category of all algebraic \(k\)-groups or in
the subcategory of those of height at most \(1\). Furthermore,
\[
  \mathfrak G_p(\pi):
  \mathfrak G_p(\mathcal L_1)
  \longrightarrow
  \mathfrak G_p(\mathcal L_2)
\]
identifies \(\mathfrak G_p(\mathcal L_2)\) with the quotient of
\(\mathfrak G_p(\mathcal L_1)\) by
\(\mathfrak G_p(\mathcal L_0)\) in the category of algebraic
\(k\)-groups.

\medskip
\noindent\textbf{Remark 8.1.4.}\footnote{The editors added this
remark, communicated by O.~Gabber, for use in \SGARefLink{sga3:VIIA:8.3.1}{8.3.1}.}
\phantomsection
\label{sga3:VIIA:8.1.4}

Let \(\varphi:G\to H\) be a morphism of \(k\)-groups and put
\(K=\ker(\varphi)\). Suppose that \(\varphi\) is covering for
the (fpqc) topology. This is in particular the case if it is
covering for a coarser topology, such as the (fppf) topology.
Then \(\varphi\) is a \(K\)-torsor over \(H\) (cf.\ IV,
\SGARefLink{sga3:iv:corollary:5:1:7:1}{5.1.7.1}). Moreover, by IV, \SGARefLink{sga3:iv:proposition:6:3:1}{6.3.1}, there exists a covering of
\(H\) by affine open subsets \(S_i\), and for each \(i\) an
affine faithfully flat morphism \(T_i\to S_i\) factoring through
\(\varphi\). Now \(G\times_HT_i\) is \(T_i\)-isomorphic to
\(K\times T_i\), hence faithfully flat over \(T_i\). By (fpqc)
descent, \(G\times_HS_i\to S_i\) is faithfully flat, and
therefore \(\varphi\) is faithfully flat.

Conversely, if \(\varphi\) is faithfully flat and quasi-compact,
respectively faithfully flat and locally of finite
presentation, then it is covering for the (fpqc), respectively
(fppf), topology (cf.\ IV, \SGARefLink{sga3:iv:proposition:6:3:1}{6.3.1}). Finally, a morphism of
sheaves is covering if and only if it is an epimorphism
(cf.\ IV, \SGARefLink{sga3:iv:proposition:4:4:3}{4.4.3}). In particular, a quasi-compact morphism of
\(k\)-groups is faithfully flat if and only if it is an
epimorphism of (fpqc) sheaves.

\medskip
\noindent\textbf{8.2. Proposition.}
\phantomsection
\label{sga3:VIIA:8.2}
\textit{Consider an exact sequence\footnote{Thus \(\pi\) is
faithfully flat and \(\tau\) identifies \(G'\) with
\(\ker\pi\), so \(G''\) represents the (fppf) quotient sheaf of
\(G\) by \(G'\); cf.\ VIA, \SGARefLink{thm:sga3-VIA-3.2}{3.2} and \SGARefLink{sga3:VIIB:5.2:theorem}{5.2}.} of algebraic groups
over a field \(k\) of characteristic \(p>0\):}
\begin{center}
  \phantomsection
  \label{sga3:VIIA:8.2-algebraic-groups-exact-sequence}
  \(\displaystyle
1 \longrightarrow G'
  \xrightarrow{\tau} G
  \xrightarrow{\pi} G''
  \longrightarrow 1 .
\)

\end{center}
\textit{Consider the following assertions:
\begin{enumerate}
  \item[\textup{(i)}] \(\pi\) is smooth;
  \item[\textup{(ii)}] \(G'\) is smooth;
  \item[\textup{(iii)}] for every integer \(n>0\), the sequence
  induced by \(\tau\) and \(\pi\),
\end{enumerate}}
\begin{center}
  \phantomsection
  \label{sga3:VIIA:8.2-frobenius-kernel-exact-sequence}
  \(\displaystyle
1 \longrightarrow \operatorname{Fr}^nG'
  \longrightarrow \operatorname{Fr}^nG
  \longrightarrow \operatorname{Fr}^nG''
  \longrightarrow 1 .
\)

\end{center}
\textit{\noindent is exact;
\begin{enumerate}
  \item[\textup{\SGAInlineRefTarget{sga3:VIIA:8.2:item:iv}(iv)}] the morphism
  \[
    \operatorname{Fr}\pi:
    \operatorname{Fr}G\longrightarrow\operatorname{Fr}G''
  \]
  is an epimorphism of (fppf) sheaves;
  \item[\textup{\SGAInlineRefTarget{sga3:VIIA:8.2:item:v}(v)}] the morphism
  \[
    \operatorname{Lie}(\pi):
    \operatorname{Lie}(G)\longrightarrow\operatorname{Lie}(G'')
  \]
  is surjective.
\end{enumerate}
Then
\[
  \textup{(i)}\Longleftrightarrow\textup{(ii)}
  \Longrightarrow\textup{(iii)}
  \Longrightarrow\textup{(iv)}
  \Longleftrightarrow\textup{(v)},
\]
and all five assertions are equivalent when \(G\) is smooth
over \(k\).}

Indeed, \SGARefLink{sga3:vii:paragraph:5-1:equation:i}{(i)} is equivalent to \SGARefLink{sga3:vii:paragraph:5-1:equation:ii}{(ii)} by VIB, \SGARefLink{prop:sga3-VIB-9.2}{9.2}~\SGARefLink{concl:sga3-VIB-9.2-vii}{(vii)}, and it is
clear that \SGARefLink{concl:sga3-VIB-9.2-iii}{(iii)} implies \SGARefLink{sga3:VIIA:8.2:item:iv}{(iv)}. The equivalence of \SGARefLink{sga3:VIIA:8.2:item:iv}{(iv)} and \SGARefLink{sga3:VIIA:8.2:item:v}{(v)}
follows from \SGARefLink{sga3:VIIA:8.1.3}{8.1.3}.

The implication \SGARefLink{concl:sga3-VIB-9.2-ii}{(ii)} \(\Rightarrow\) \SGARefLink{concl:sga3-VIB-9.2-iii}{(iii)} follows from the
diagram
\begin{center}
  \phantomsection
  \label{sga3:VIIA:8.2-frobenius-exact-rows-diagram}
  \resizebox{0.98\linewidth}{!}{%
\begin{tikzcd}[ampersand replacement=\&,row sep=3.6em,column sep=3.5em]
1 \arrow[r]
  \& G' \arrow[r,"\tau"] \arrow[d,swap,"\operatorname{Fr}^n(G'/k)"]
  \& G \arrow[r,"\pi"] \arrow[d,swap,"\operatorname{Fr}^n(G/k)"]
  \& G'' \arrow[r] \arrow[d,swap,"\operatorname{Fr}^n(G''/k)"]
  \& 1 \\
1 \arrow[r]
  \& G'^{(p^n)} \arrow[r,"\tau^{(p^n)}"']
  \& G^{(p^n)} \arrow[r,"\pi^{(p^n)}"']
  \& G''^{(p^n)} \arrow[r]
  \& 1
\end{tikzcd}%
}

\end{center}
whose two rows are exact. Since
\(\operatorname{Fr}^n(G'/k)\) is an epimorphism of (fppf)
sheaves by \SGARefLink{sga3:VIIA:8.3.1}{Corollary 8.3.1} below, \(\pi\) induces an
epimorphism from \(\operatorname{Fr}^nG\) onto
\(\operatorname{Fr}^nG''\); here one uses the snake lemma for
sheaves of groups, not necessarily commutative.

Likewise, suppose that \(G\) is smooth over \(k\). Then
\(\operatorname{Fr}(G/k)\) is an epimorphism. If
\(\operatorname{Fr}\pi\) is also an epimorphism, the same snake
lemma, applied to the preceding diagram with \(n=1\), shows
that \(\operatorname{Fr}(G'/k)\) is an epimorphism. Hence \(G'\)
is smooth over \(k\) by \SGARefLink{sga3:VIIA:8.3.1}{8.3.1} below.

\medskip
\noindent\textbf{8.3. Proposition.}\footnote{The editors
replaced ``algebraic'' by ``locally of finite type.''}
\phantomsection
\label{sga3:VIIA:8.3}
\textit{Let \(G\) be a group locally of finite type over a field
\(k\) of characteristic \(p>0\). There exists an integer
\(n_0\) such that
\[
  G/(\operatorname{Fr}^nG)
\]
is smooth over \(k\) for every \(n\ge n_0\).}

The construction of \(G/(\operatorname{Fr}^nG)\) commutes with
extension of the base field (\SGARefLink{sga3:VIIA:4.1.1}{4.1.1} and VIA, \SGARefLink{thm:sga3-VIA-3.3.2}{3.3.2}), so we may
suppose that \(k\) is perfect. Then \(G_{\mathrm{red}}\) is a
\(k\)-group locally of finite type (cf.\ VIA, \SGARefLink{sga3:VIA:0.2}{0.2}), and, putting
\[
  H=G_{\mathrm{red}}\backslash G,
\]
we have the following commutative exact diagram:
\begin{center}
  \phantomsection
  \label{sga3:VIIA:8.3-reduced-quotient-frobenius-diagram}
  \resizebox{0.94\linewidth}{!}{%
\begin{tikzcd}[ampersand replacement=\&,row sep=3.8em,column sep=3.7em]
1 \arrow[r]
  \& G_{\mathrm{red}} \arrow[r] \arrow[d,"\operatorname{Fr}^n(G_{\mathrm{red}}/k)"]
  \& G \arrow[r] \arrow[d,"\operatorname{Fr}^n(G/k)"]
  \& H \arrow[d,"\operatorname{Fr}^n(H/k)"] \\
1 \arrow[r]
  \& G_{\mathrm{red}}^{(p^n)} \arrow[r]
  \& G^{(p^n)} \arrow[r]
  \& H^{(p^n)}\, .
\end{tikzcd}%
}

\end{center}

The scheme \(H\) is the spectrum of a finite local \(k\)-algebra
with residue field \(k\) (cf.\ VIA, \SGARefLink{prop:sga3-VIA-5.6.1}{5.6.1}). There is therefore
an integer \(n_0\) such that, for every \(n\ge n_0\),
\(\operatorname{Fr}^n(H/k)\) factors through the unit section of
\(H^{(p^n)}\). It follows that, for \(n\ge n_0\),
\(\operatorname{Fr}^n(G/k)\) factors through
\(G_{\mathrm{red}}^{(p^n)}\), and VIA, \SGARefLink{prop:sga3-VIA-5.4.1}{5.4.1} gives the
commutative diagram
\begin{center}
  \phantomsection
  \label{sga3:VIIA:8.3-frobenius-factorization-triangle}
  \begin{tikzcd}[row sep=3.6em,column sep=4.6em]
G \arrow[r,"\operatorname{Fr}^n(G/k)"] \arrow[d,swap,"\pi"]
  & G_{\mathrm{red}}^{(p^n)} \\
G/(\operatorname{Fr}^nG) \arrow[ur,swap,"i"]
\end{tikzcd}

\end{center}
where \(i\) is a closed immersion and \(\pi\) is the canonical
projection. Moreover, \(i\) induces a homeomorphism of the
underlying topological spaces, so it is an isomorphism. Since
\(k\) is perfect, \(G_{\mathrm{red}}^{(p^n)}\) is smooth over
\(k\) (VIA, \SGARefLink{prop:sga3-VIA-1.3.1}{1.3.1}). Thus \(G/(\operatorname{Fr}^nG)\) is smooth
over \(k\) for every \(n\ge n_0\).

\medskip
\noindent\textbf{8.3.1. Corollary.}
\phantomsection
\label{sga3:VIIA:8.3.1}
\textit{Let \(G\) be a group locally of finite type over a field
\(k\) of characteristic \(p>0\), and let \(n\ge1\).\footnote{The
editors made explicit the equivalence of \SGARefLink{sga3:vii:paragraph:5-1:equation:ii}{(ii)} and \SGARefLink{sga3:VIIA:8.3.1:item:iii}{(iii)} and
expanded the proof.} The following conditions are equivalent:
\begin{enumerate}
  \item[\textup{(i)}] \(G\) is smooth over \(k\);
  \item[\textup{(ii)}]
  \(\operatorname{Fr}^n(G/k):G\to G^{(p^n)}\) is an
  epimorphism of (fppf) sheaves;
  \item[\textup{\SGAInlineRefTarget{sga3:VIIA:8.3.1:item:iii}(iii)}]
  \(\operatorname{Fr}^n(G/k):G\to G^{(p^n)}\) is faithfully
  flat.
\end{enumerate}}

Since \(G\) is locally of finite type over \(k\),
\(\operatorname{Fr}^n(G/k)\) is of finite presentation, so the
equivalence of \SGARefLink{sga3:vii:paragraph:5-1:equation:ii}{(ii)} and \SGARefLink{sga3:VIIA:8.3.1:item:iii}{(iii)} follows from \SGARefLink{sga3:VIIA:8.1.4}{8.1.4}. Suppose that
\(G\) is smooth over \(k\), and hence reduced. Since
\(\operatorname{Fr}^n(G/k)\) is surjective, it is faithfully
flat (cf.\ VIA, \SGARefLink{cor:sga3-VIA-6.2}{6.2} or VIB, \SGARefLink{prop:sga3-VIB-1.3}{1.3}).

Conversely, suppose that \(\operatorname{Fr}^n(G/k)\) is
faithfully flat. The morphism
\[
  \operatorname{Fr}^n(G^{(p^n)}/k)
\]
is obtained from \(\operatorname{Fr}^n(G/k)\) by base change
(cf.\ \SGARefLink{sga3:VIIA:4.1.3}{4.1.3}), so it too is faithfully flat, as is the composite
\[
  \operatorname{Fr}^{2n}(G/k):
  G\longrightarrow G^{(p^n)}
  \longrightarrow G^{(p^{2n})}.
\]
It follows that, for every \(m\in\mathbb N\),
\[
  \operatorname{Fr}^{mn}(G/k):
  G\longrightarrow G^{(p^{mn})}
\]
is faithfully flat. Hence it induces an isomorphism
\[
  G/(\operatorname{Fr}^{mn}G)
  \simeq
  G^{(p^{mn})}
\]
(cf.\ VIA, \SGARefLink{prop:sga3-VIA-5.4.1}{5.4.1}). By \SGARefLink{sga3:VIIA:8.3}{Proposition 8.3},
\(G^{(p^{mn})}\) is smooth over \(k\) for \(m\) sufficiently
large, and therefore \(G\) is smooth as well by (fpqc) descent
(cf.\ EGA IV\(_4\), 17.7.1).

